
\SGASection{1}{Definition of unipotent algebraic groups}
\SGARefTarget{sga3:xvii:section:1}

\SGAStatement{Definition 1.1}{%
An algebraic group \(G\) defined over an algebraically closed field \(k\) is
called \emph{unipotent} if \(G\) admits a composition series whose successive
quotients are isomorphic to algebraic subgroups of \((\mathbb G_a)_k\).}
\SGARefTarget{sga3:xvii:definition:1:1}

\SGAStatement{Proposition 1.2}{%
Let \(k\) be an algebraically closed field, let \(K\) be an algebraically closed
extension of \(k\), and let \(G\) be an algebraic group defined over \(k\).
Then \(G\) is unipotent if and only if \(G_K\) is unipotent.}
\SGARefTarget{sga3:xvii:proposition:1:2}

The necessity is clear, since a composition series remains a composition series
after extension of the base.

The proof of sufficiency is standard.  The field \(K\) is the filtered direct
limit of its finitely generated \(k\)-subalgebras.  By Exposé
VI\textsubscript{B}, \S10, one can find a finitely generated \(k\)-subalgebra
\(A\) of \(K\), a composition series \(G_i\) of \(G_S\), where
\(S=\operatorname{Spec}A\), and immersions
\[
  u_i:H_i=G_i/G_{i+1}\longrightarrow(\mathbb G_a)_S .
\]
To prove that \(G\) is unipotent, it therefore suffices to make a \(k\)-valued
base change \(A\to k\).  This is possible because
\(\operatorname{Hom}_{k\text{-alg}}(A,k)\) is nonempty: \(A\) is a nonzero
finitely generated algebra over the algebraically closed field \(k\).

\SGAStatement{Definition 1.3}{%
Let \(G\) be an algebraic group defined over a field \(k\).  We say that \(G\)
is \emph{unipotent} if there exists an algebraically closed extension
\(\bar{k}\) of \(k\) such that \(G_{\bar{k}}\) is unipotent in the sense of
\SGARefLink{sga3:xvii:definition:1:1}{Definition 1.1}.}
\SGARefTarget{sga3:xvii:definition:1:3}

By \SGARefLink{sga3:xvii:proposition:1:2}{Proposition 1.2}, this property is independent of the chosen algebraically
closed extension \(\bar{k}\).

\SGAStatement{Definition 1.4}{%
Let \(G\) be an algebraic group defined over a field \(k\), and let \(H\) be an
algebraic group defined over an extension \(k'\) of \(k\).  We say that \(H\)
is a \emph{form} of \(G\) over \(k'\) if the algebraic groups
\(G_{\bar{k}'}\) and \(H_{\bar{k}'}\) become isomorphic over an algebraic
closure \(\bar{k}'\) of \(k'\).  As above, this property is independent of
the choice of \(\bar{k}'\).  We shall also say that \(H\) is a
\emph{twisted form} of \(G\).}
\SGARefTarget{sga3:xvii:definition:1:4}

We can now describe the algebraic subgroups of \(\mathbb G_a\).

\SGAStatement{Proposition 1.5}{%
Let \(k\) be a field of characteristic \(p\geq0\).  An algebraic subgroup
\(H\) of \((\mathbb G_a)_k\) is of one of the following types:
\begin{enumerate}
\item[\rm(i)] \(H=0\);
\item[\rm(ii)] \(H=\mathbb G_a\);
\item[\rm(iii)]\SGAInlineRefTarget{sga3:xvii:proposition:1:5:item:iii} if \(p>0\), \(H\) is an extension of a twisted constant group
  \((\mathbb Z/p\mathbb Z)^r\) by a radicial group
  \(\boldsymbol\alpha_{p^n}\), where \(r,n\geq0\) and \(r+n>0\).
  If, moreover, \(k\) is perfect, this extension is necessarily trivial.
\end{enumerate}}
\SGARefTarget{sga3:xvii:proposition:1:5}
